\documentclass[twocolumn,10pt,a4paper]{article}
\usepackage[margin=0.75in]{geometry}
\usepackage[utf8]{inputenc}
\usepackage[T1]{fontenc}
\usepackage{lmodern}
\usepackage{microtype}
\usepackage{cite}
\usepackage{url}
\usepackage{hyperref}
\usepackage{xcolor}
\usepackage{enumitem}
\usepackage{tabularx}
\usepackage{booktabs}
\hypersetup{colorlinks=true,citecolor=blue,linkcolor=blue,urlcolor=blue}

\title{A Fixture-Driven Research Harness for Falsifiable Claims in Agentic O-RAN Automation}

\author{David Kypuros \\
\small Red Hat \\
\small \texttt{dkypuros@redhat.com}}

\date{31 May 2026}

\begin{document}
\maketitle

\begin{abstract}
Agentic-AI-for-networks papers frequently make architectural claims that are difficult to falsify because the implementation is proprietary, partially specified, or not runnable independently. We describe the research engineering methodology behind a closed-loop Day-2 remediation harness for O-RAN Cloud RAN that takes the opposite posture: every architectural claim anchors to a file, every file carries a citation header, and a 10-check verify gate produces a binary pass/fail any reviewer can reproduce on a clean clone in under five seconds. The methodology has four anchoring layers (citation headers, JSON Schema contracts, a bidirectional conformance index, and end-to-end fixture replay) and an explicit stub-boundary discipline that names which components carry the paper's claims and which exist only to isolate the real components from external dependencies. The LLM-neutral design is enforced structurally rather than by convention, which preserves determinism in the verify gate. We argue this style of research substrate produces falsifiable outputs independent of LLM provider, API availability, or model-version drift.
\end{abstract}

\section{Introduction}

For a software harness to serve as a research substrate (as opposed to a product prototype), it should produce claims that a reviewer can falsify mechanically. Three properties matter: claims are anchored to specific files, the substrate runs reproducibly without external dependencies, and stub boundaries are declared rather than hidden. The dominant pattern in agentic-AI-for-networks publications fails one or more of these properties. Implementations are described in prose without file-level anchors; the runtime depends on a live cloud, a live cluster, or a particular API key; and the stub boundary (the line between code under test and the systems being mocked) is implicit.

A prior multivendor diagnostic publication \cite{multivendor_blog_2026} established a fault-localization substrate for Cloud RAN troubleshooting. The closed-loop extension we describe in companion papers takes the same posture: the research artifact is the runnable open-source repository, the architectural narrative is anchored at file plus line, and the verify gate is the falsification mechanism. This paper documents the methodology by which the three substrate properties are achieved. The repository under discussion is open-source under Apache-2.0 at \url{https://github.com/dkypuros/oran-agent-harness}.

The remainder of the paper is organized as follows. Section \ref{sec:anchoring} documents the four-layer claim anchoring system. Section \ref{sec:gate} describes the verify gate's ten checks and the failure mode each one guards. Section \ref{sec:stub} describes the stub-boundary discipline. Section \ref{sec:llm} describes the LLM-neutral design as a research methodology requirement, not only a production concern. Section \ref{sec:compare} compares the approach to adjacent research engineering practices. Section \ref{sec:repro} states the reproducibility entry point. Section \ref{sec:conclude} concludes.

\section{The Claim Anchoring System}
\label{sec:anchoring}

Four layers of anchoring carry the load.

\subsection{Layer 1: Citation headers on every stub}

Every YAML and JSON file authored under \texttt{harness/} and \texttt{scenarios/} opens with a \texttt{\_conforms\_to} or equivalent block naming the spec, the spec section, the spec version, and the bibliography reference numbers. The verify gate's first check enforces this: any new file added without a citation header fails the gate. The mechanism turns documentation into a contract. Adding a routing rule, a schema, a guardrail, or a scenario fixture without naming the spec it conforms to is a build failure, not a documentation oversight.

\subsection{Layer 2: JSON Schema draft-07 contracts}

Five schemas define the shapes the pipeline produces and consumes: \texttt{FaultPayload.json}, \texttt{RCA.json}, \texttt{RemediationProposal.json}, \texttt{AuditEvent.json}, and \texttt{ReversibilityProfile.json}. Each scenario's \texttt{fault\_payload.json} validates against \texttt{FaultPayload.json}; each scenario's \texttt{audit\_event.json} validates against \texttt{AuditEvent.json} with the embedded reversibility profile validating against its own schema. The gate's ninth check runs the validation. A claim like "the audit event always carries a populated reversibility profile" stops being a prose assertion and becomes a schema validation outcome.

\subsection{Layer 3: Bidirectional conformance index}

A single file \texttt{harness/conformance.md} is a two-column table mapping every authored artifact in the harness tree to its spec lineage and bibliography references. The gate's fourth check confirms both directions: every file under \texttt{harness/} or \texttt{scenarios/} appears in the table, and every row in the table points at a file that exists. The mechanism prevents the index from rotting silently as the codebase changes. A new file without an entry fails the gate; an entry pointing at a deleted file fails the gate.

\subsection{Layer 4: End-to-end fixture replay}

The gate's tenth check runs the deterministic walker as a subprocess against all five committed scenario fault payloads. For each scenario, the actual audit event output is compared against the committed expected audit event on seven material fields: \texttt{targetLayer}, \texttt{taxonomyMatch}, \texttt{actionType}, \texttt{ocloudInternalPath}, \texttt{actionPayloadRef}, \texttt{dryRun}, \texttt{requiresHumanApproval}, plus the \texttt{reversibility\_profile.confidence\_in\_reversibility} field. Any divergence on any material field fails the gate. The mechanism turns architectural claims about routing decisions, action types, and the seven-field reversibility contract into mechanically falsifiable assertions.

\section{The Verify Gate}
\label{sec:gate}

The gate is a single Python script at \texttt{scripts/verify.py}. It runs ten checks in sequence and exits 0 on 10 of 10 pass, non-zero with specific failure rows otherwise. Each check guards a specific failure mode.

\textbf{Check 1, citation\_headers.} Guards against undocumented stubs. A new YAML or JSON file added without a \texttt{\_conforms\_to} block fails the check.

\textbf{Check 2, json\_parse.} Guards against syntax rot in any JSON file across the public tree.

\textbf{Check 3, yaml\_parse.} Same for YAML files.

\textbf{Check 4, conformance\_complete.} Bidirectional index integrity (Layer 3 above).

\textbf{Check 5, em\_dashes.} Guards against a specific authorship tell that the author has committed to eliminating from published copy. The check reports zero permitted occurrences of the em dash codepoint (U+2014) in the public tree.

\textbf{Check 6, no\_leakage.} Guards against accidentally committing \texttt{.env} files, \texttt{.local/} content, credentials, or large local demo artifacts. Runs against \texttt{git ls-files} so it catches anything the working tree might hide.

\textbf{Check 7, readme\_structure.} Guards against the README drifting from the three stated repository goals. Enforces line-count bounds and the presence of the three-goal lede.

\textbf{Check 8, file\_counts.} Guards against silent file additions or deletions that change the harness surface without updating the index. The count is asserted at known totals.

\textbf{Check 9, schema\_validation.} Scenario data validates against the JSON Schemas (Layer 2 above). Auto-skips if the \texttt{jsonschema} package is not installed, with a warning.

\textbf{Check 10, walker\_e2e.} The falsifiability check (Layer 4 above). Spawns the walker as a subprocess, captures the audit event, compares to the committed expected output on material fields.

Run-time of the entire gate against the reference repository state is approximately five seconds on a developer-grade workstation. The dependency surface is Python 3.8 or later, PyYAML, and (optionally) jsonschema. No cloud account, no live cluster, no LLM API key is required for the deterministic gate.

\section{Stub Boundary Discipline}
\label{sec:stub}

The harness is explicit about what is real and what is stubbed. The methodology argument is that the components carrying the paper's claims must be real; the components that exist only to isolate them from external dependencies should be stubbed cleanly.

\textbf{Real components.} The \texttt{harness/runtime/router.py} module performs the deterministic taxonomy lookup against the authored \texttt{harness/taxonomy.yaml}. It is exercised end-to-end in check 10 of the verify gate. The \texttt{harness/runtime/guardrail.py} module evaluates the five-gate guardrail against the authored \texttt{harness/guardrails.yaml}. Both modules' deterministic behavior is the source of the paper's claims about routing and gating. They are real Python code executing real YAML inputs.

The repository also ships a real Python function call that crosses the O-RAN O2 IMS edge from the architecture diagram: \texttt{5G\_O-RAN\_SIM/oam/o2ims\_stub.py}'s \texttt{deploy\_request} function is invoked by the router on every infrastructure-layer route, and the response envelope is attached to the audit. The stub is shape-only (no network call, no real reconcile loop), but the function call is real, and the test \texttt{test\_o2ims\_dispatch\_attached\_on\_down\_route} asserts that the response shape is contract-correct.

\textbf{Stubbed components.} The Agentic Gateway, the four MCP servers (\texttt{mcp-platform}, \texttt{mcp-ran}, \texttt{mcp-hardware}, \texttt{mcp-ocloud}), the three domain agents (Platform, RAN, Hardware), and the Digital Twin substrate are stubbed. The walker's module docstring names them directly. The stub boundary is declared, not hidden. Each scenario's canned values (RCA timestamp, contributing signals, action target, blast radius, sandbox verdict, SMO dispatch outcome, O2 IMS dispatch outcome) live in a single file at \texttt{harness/runtime/scenario\_stubs.json}, which serves as the explicit single source of truth for stubbed data.

This design decision matters for the research substrate property. The real components (router, guardrail, o2ims function call) are the ones that carry the paper's architectural claims, and they are exercised against authored YAML and JSON inputs the verify gate validates. The stub components are wrappers that isolate the real components from external dependencies (a live MCP server, a live LLM endpoint, a live Twin cluster) so the gate can run without network access, credentials, or vendor-specific tooling. A reviewer can change the stubbed inputs and observe the real components' deterministic response.

\section{The LLM-Neutral Substrate as a Research Property}
\label{sec:llm}

The harness's third declared contribution is an LLM-neutral substrate. The motivation is partly production-pragmatic (avoid lock-in to a specific provider) but the research methodology value is stronger.

If the runtime embedded a specific LLM provider's SDK or a specific model identifier, the verify gate's tenth check (walker end-to-end) would have nondeterministic output. Different reviewers running the gate against the same repository state at the same git SHA would see different outcomes depending on provider availability, model version, or even response-time jitter. The research substrate property would collapse.

The harness enforces LLM neutrality structurally rather than by convention. The routing rule's ambiguous path raises \texttt{NotImplementedError} in the default code path, preserving determinism in the verify gate. The live LLM path is gated behind an environment variable (\texttt{ORAN\_LLM\_MODE=live}) and a separate test (\texttt{test\_route\_ambiguous\_with\_llm\_mode\_live}) that skips by default to keep the gate result independent. The structural enforcement means a reviewer cannot accidentally make the gate provider-dependent.

The LLM is therefore not load-bearing on the routing decision the verify gate proves. This is the right division of labor for a research substrate: the deterministic claims are reproducible at compile time, the LLM-assisted claims are gated behind a separate execution mode that does not interfere with the deterministic gate. A v1 deployment that wants real LLM behavior on the ambiguous path can light up the env-gated seam without touching the deterministic core.

\section{Comparison to Adjacent Research Practices}
\label{sec:compare}

Three adjacent research engineering practices illustrate the niche the harness occupies.

\textbf{Unit tests without spec-anchored fixtures.} Most research code ships with tests that confirm internal consistency (a function returns what its signature promises) but not external spec conformance. The harness's check 10 binds the test output to committed expected audit fixtures, and those fixtures themselves conform to spec-cited JSON Schemas. A change that breaks the architectural claim breaks the gate.

\textbf{Integration tests against live services.} Some research projects integration-test against a live cluster or API. This is honest about end-to-end behavior but defeats reproducibility: a reviewer without the live service cannot rerun the test. The harness preserves end-to-end coverage through fixture replay against committed canned inputs, which a reviewer can re-execute on a clean clone.

\textbf{Pure simulation.} Some research uses simulators that never touch real code paths through standardized interfaces. The harness occupies a middle position: the deterministic routing and guardrail logic are real Python executing real YAML; the simulated component is the partner system (the O-Cloud Manager, the partner SMO, the BMC) which is stubbed at the shape-correct envelope boundary. A v1 deployment can swap the stubs for live targets behind the same envelope.

\section{Reproducibility Entry Point}
\label{sec:repro}

A reader who wants to reproduce the harness verify gate runs three commands.

\begin{verbatim}
git clone \
  https://github.com/dkypuros/oran-agent-harness
pip install pyyaml jsonschema
python3 scripts/verify.py
\end{verbatim}

That is the complete reproducibility claim for the deterministic gate. No cloud account, no live cluster, no LLM API key required. The reference state at the time of writing is ten of ten checks pass.

Beyond the gate, the repository ships two lab postures. The macOS posture at \texttt{macbook\_lab/} uses Docker Desktop and a multi-container compose file to stand up a fuller stack with HTTP wrappers around the platform stubs. The Linux posture at \texttt{linux\_lab/} uses podman 5.4.1 on Fedora 42 with the same compose file. A verified run document at \texttt{linux\_lab/verified\_runs/2026-05-31\_fedora42\_podman.md} captures the live cross-runtime verification.

\section{Honest Scope}

The methodology produces falsifiable claims for the bounded set of behaviors the verify gate covers. Several explicit limits.

\textbf{Fixtures are canned.} Scenario inputs and the partner-system stubbed responses do not come from a running O-Cloud or a real linuxptp daemon. A real cluster could produce inputs that fail validation against the JSON Schemas; that would be informative but is out of scope for the gate.

\textbf{Schemas are harness-authored.} The JSON Schema draft-07 contracts at \texttt{harness/schemas/} are the harness's own approximations of the spec envelopes. The schemas are not published by O-RAN, TM Forum, or any other standards body. A real spec-conformance test would use the standards bodies' own test suites (for example the O-RAN SC compliance test suite \cite{oran_sc_it_test}); that is a v1 work item.

\textbf{Scenario coverage is partial.} The five committed scenarios cover the primary delivery path branches (MCO, KMM, Metal3, dual-route accepted, dual-route SMO-rejected) but 15 of the 20 taxonomy entries have no walkthrough fixture. A future research artifact could claim broader coverage once those fixtures exist.

\textbf{Some signals are stubbed values.} EvalOps confidence telemetry, Agentic Recoverability validation history, the full Digital Twin convergence envelope, and the live LLM response on the ambiguous path are described in the architectural narrative and the schemas but populated as stubbed values in \texttt{scenario\_stubs.json} for v0. Their real-time operation is a v1 deliverable.

\section{Conclusion}
\label{sec:conclude}

This paper has described the research engineering methodology behind a closed-loop Day-2 remediation harness for O-RAN Cloud RAN. Four anchoring layers (citation headers, JSON Schema contracts, a bidirectional conformance index, fixture replay) and a 10-check verify gate make every architectural claim mechanically falsifiable. The stub boundary is declared, not hidden, and the components carrying the paper's claims are real code exercised against authored inputs. LLM neutrality is enforced structurally, which preserves verify-gate determinism. The combination produces a research substrate that any reviewer can reproduce on a clean clone in under five seconds. We argued this is a meaningful improvement over agentic-AI publications that ship prose architectural claims without spec-anchored falsification mechanisms. v0 limits are declared; the v1 path is named.

\section*{Acknowledgements}

The methodology builds on a prior multivendor diagnostic substrate co-authored with S. Kiani Mehr, N. Korati Prasanna, M. Vazquez Cebrian, S. Srikanthan, and P. Venkatesh \cite{multivendor_blog_2026}. Source code is Apache-2.0 licensed at \url{https://github.com/dkypuros/oran-agent-harness}.

\bibliographystyle{plain}
\bibliography{refs}

\end{document}
